\documentclass[10pt,a4paper]{article}

\usepackage{iftex}
\ifPDFTeX
  \usepackage[T2A]{fontenc}
  \usepackage[utf8]{inputenc}
  \usepackage[russian,english]{babel}
\else
  \usepackage{fontspec}
  \usepackage[russian,english]{babel}
  \setmainfont{DejaVu Serif}
  \setsansfont{DejaVu Sans}
  \setmonofont{DejaVu Sans Mono}
\fi
\usepackage{amsmath,amssymb}
\usepackage{array,longtable}
\usepackage{enumitem}
\usepackage{geometry}
\usepackage{hyperref}
\usepackage{url}

\geometry{left=18mm,right=18mm,top=16mm,bottom=18mm}
\setlength{\parindent}{3.5mm}
\setlength{\parskip}{0.25em}
\setlist{nosep,leftmargin=5mm}
\renewcommand{\arraystretch}{1.18}
\sloppy

\newcommand{\code}[1]{\ifmmode\text{\ttfamily\detokenize{#1}}\else\texttt{\detokenize{#1}}\fi}
\newcommand{\A}[1]{\mathcal{A}_{#1}}
\newcommand{\Contract}[1]{\mathcal{C}_{#1}}
\newcommand{\Env}{\mathcal{A}_{ENV}}
\pdfstringdefDisableCommands{\def\A#1{A_#1}\def\Contract#1{C_#1}}

\title{Формализация Application / Service Layer в иерархической SDN-управляемой 6G ISAC-сети}
\author{Техническая документация уровня}
\date{8 июня 2026 г.}

\begin{document}

\maketitle

\section*{Аннотация}

Документ задает формальную спецификацию Application / Service Layer (прикладного или сервисного уровня) для SDN-управляемой 6G ISAC-сети. Этот уровень формирует требования к сервису, SLA (service-level agreement, соглашение об уровне обслуживания), критичность и тип нагрузки, но не управляет радиоресурсами напрямую. Он получает от SDN/RIC результат допуска, деградации или отказа и контролирует выполнение SLA в конечных классах. Модель не содержит ИИ и не предполагает автоматического обучения; сервисный уровень описывается как аппроксимированный контрактный timed automaton.

\section{Назначение и границы уровня}

Application / Service Layer задает требования к сети через northbound interface (северный интерфейс) SDN-контроллера. Такое разделение приложений и уровня управления соответствует SDN-архитектуре сетей 5G/6G \cite{mosudi_sdn,zolanvari_sdn}. В ISAC-сети сервис может требовать не только communication QoS (качество передачи данных), но и sensing QoS (качество радиозондирования): период обновления, допустимую свежесть, точность, покрытие зоны и минимальную способность обнаружения малоразмерных БПЛА.

Граница уровня:
\begin{itemize}
  \item Сервис формирует request (запрос), criticality (критичность), SLA constraints (ограничения SLA) и demand type (тип требования).
  \item Сервис не выдает команд \code{beam_update_cmd!}, \code{sensing_boost_cmd!}, \code{flow_mod!} или команд распределения физических ресурсов.
  \item Сервис получает \code{service_accept!}, \code{service_degraded!} или \code{service_reject!} от SDN/RIC.
  \item Сервис контролирует SLA только в конечных классах, пригодных для model checking (формальной проверки модели).
\end{itemize}

Если сервисный уровень напрямую управляет beam management (управлением лучом), PRS-pattern (паттерном позиционно-измерительных сигналов) или расписанием OFDM-ресурсов, нарушается разделение уровней. Такие действия должны выполняться через SDN/RIC и MAC.

\section{Архитектурное положение сервисного уровня}

Полная композиция:
\[
\A{TOTAL}=\A{SVC}\parallel \A{SDN}\parallel \A{MAC}\parallel \A{PHY}\parallel \Env.
\]

Сервисный уровень является верхним контрактным компонентом:
\[
Gar_{SVC}\Rightarrow Ass_{SDN},\qquad
Gar_{SDN}\Rightarrow Ass_{SVC}.
\]

Для локальной проверки:
\[
\A{SYS}^{SVC}=\A{SVC}\parallel \A{ENV\_SVC},
\qquad
\A{ENV\_SVC}=\pi_{SVC}(\Env).
\]

\(\A{ENV\_SVC}\) моделирует поступление новых сервисных требований, изменение критичности, завершение сессии и наблюдаемое нарушение SLA. Она не должна противоречить SDN-проекции общей среды: если SDN вернул \code{service_reject!}, сервис не может одновременно перейти в \code{Accepted} без явного degraded-решения.

\section{Базовая модель контрактного автомата}

Компонент сервисного уровня задается как:
\[
\A{i}=(L_i,l_i^0,C_i,V_i,E_i,\Inv_i,In_i,Out_i,Ass_i,Gar_i).
\]

Переход:
\[
e=(l,g,a,r,u,l').
\]

Модель проверяет своевременность реакции на решения SDN/RIC и нарушения SLA в семантике timed automata \cite{alur_dill}. Она не доказывает истинную вероятность обнаружения БПЛА или истинную надежность канала; эти параметры должны поступать в виде конечных классов от нижних уровней и внешнего estimator (измерительного или аналитического оценочного слоя).

\section{Дискретная абстракция сервиса}

\[
\alpha_{SVC}:X_{SVC}^{req}\times X_{SVC}^{obs}\rightarrow X_{SVC}^{disc}.
\]

\(X_{SVC}^{req}\) содержит latency bound (границу задержки), reliability target (целевую надежность), throughput demand (потребность в пропускной способности), sensing update period (период обновления sensing), sensing accuracy class (класс точности радиозондирования), coverage requirement (требование покрытия). \(X_{SVC}^{obs}\) содержит наблюдаемый результат: принятие, отказ, деградированный режим, SLA warning (предупреждение SLA), SLA violation (нарушение SLA).

\begin{longtable}{p{0.24\linewidth}p{0.38\linewidth}p{0.30\linewidth}}
\hline
\textbf{Переменная} & \textbf{Значения} & \textbf{Смысл}\\
\hline
\code{ServiceClass} & \code{SVC_URLLC}, \code{SVC_EMBB}, \code{SVC_MMTC}, \code{SVC_ISAC_SENSING}, \code{SVC_UAV}, \code{SVC_INDUSTRIAL} & Класс сервиса.\\
\code{CriticalityClass} & \code{CRIT_LOW}, \code{CRIT_MED}, \code{CRIT_HIGH}, \code{CRIT_SAFETY} & Критичность.\\
\code{SLAClass} & \code{SLA_OK}, \code{SLA_WARN}, \code{SLA_VIOLATED} & Состояние SLA.\\
\code{AdmissionClass} & \code{ADM_PENDING}, \code{ADM_ACCEPTED}, \code{ADM_REJECTED}, \code{ADM_DEGRADED} & Решение допуска.\\
\code{DemandClass} & \code{DEMAND_COMM}, \code{DEMAND_SENS}, \code{DEMAND_JOINT} & Тип требования.\\
\code{FreshnessClass} & \code{FRESH}, \code{STALE}, \code{EXPIRED} & Давность sensing-информации.\\
\code{DetectionFreshnessClass} & \code{DET_FRESH}, \code{DET_STALE}, \code{DET_EXPIRED} & Актуальность detection-события для текущей цели.\\
\code{UpdatePeriodClass} & \code{UPD_OK}, \code{UPD_SLOW}, \code{UPD_VIOLATED} & Период обновления sensing-оценок.\\
\code{FalseAlarmClass} & \code{FA_OK}, \code{FA_HIGH}, \code{FA_CRITICAL} & Класс риска ложной тревоги.\\
\code{MissedDetectionClass} & \code{MISS_OK}, \code{MISS_HIGH}, \code{MISS_CRITICAL} & Класс риска пропуска цели.\\
\code{AccuracyClass} & \code{OK}, \code{LIMITED}, \code{FAILED} & Класс точности оценки положения цели.\\
\code{CoverageClass} & \code{OK}, \code{LIMITED}, \code{FAILED} & Класс покрытия зоны наблюдения, согласованный с PHY/SDN-интерфейсом.\\
\code{DegradedReason} & \code{DEG_NONE}, \code{DEG_STALE}, \code{DEG_LOW_PD}, \code{DEG_HIGH_PFA}, \code{DEG_UPDATE_SLOW}, \code{DEG_RESOURCE_LIMITED} & Причина деградированного сервиса.\\
\hline
\end{longtable}

\section{Композиция автоматов сервиса}

\[
\A{SVC}=\A{REQ}\parallel \A{SLA}\parallel \A{CRIT}\parallel \A{SVC\_AGG}.
\]

\begin{longtable}{p{0.18\linewidth}p{0.39\linewidth}p{0.35\linewidth}}
\hline
\textbf{Автомат} & \textbf{Основные состояния} & \textbf{Ответственность}\\
\hline
\(\A{REQ}\) & \code{ServiceIdle}, \code{RequestBuild}, \code{RequestPending}, \code{Accepted}, \code{AcceptedDegraded}, \code{Rejected}, \code{Completed} & Жизненный цикл сервисного запроса.\\
\(\A{SLA}\) & \code{SLAOk}, \code{SLAWarning}, \code{SLAViolated}, \code{SLAReported} & Контроль выполнения SLA.\\
\(\A{CRIT}\) & \code{NonCritical}, \code{Critical}, \code{SafetyCritical} & Классификация критичности.\\
\(\A{SVC\_AGG}\) & \code{DemandBuild}, \code{DemandSent}, \code{AwaitDecision}, \code{DecisionReceived} & Передача требований в SDN/RIC.\\
\hline
\end{longtable}

\section{Политики сервисного уровня}

Политики сервисного уровня не являются политиками радиоресурса. Они задают правила формирования требований, деградации и реакции на SLA.

\begin{longtable}{p{0.24\linewidth}p{0.36\linewidth}p{0.32\linewidth}}
\hline
\textbf{Политика} & \textbf{Условие} & \textbf{Результат}\\
\hline
\code{P_ADMISSION_REQUEST} & Новый сервисный спрос и корректный конечный класс требований. & Формирование \code{service_request!}.\\
\code{P_ACCEPT_DEGRADED} & SDN/RIC вернул \code{service_degraded!}, а сервис допускает деградацию. & Переход в \code{AcceptedDegraded} с контролем SLA.\\
\code{P_REJECT_REASON} & SDN/RIC вернул отказ или admission timeout. & Переход в \code{Rejected} с явной причиной.\\
\code{P_SLA_WARNING} & \code{SLAClass=SLA_WARN}. & Отчет \code{sla_warning_report!} до \(D_{sla\_warn}\).\\
\code{P_SLA_VIOLATION} & \code{SLAClass=SLA_VIOLATED}. & Отчет \code{sla_violation_report!} до \(D_{sla\_violation}\).\\
\code{P_SAFETY_CRITICAL} & \code{CRIT_SAFETY} и нарушение sensing/communication SLA. & Запрос реконфигурации, завершение сервиса или аварийный отчет.\\
\hline
\end{longtable}

Для темы обнаружения малоразмерных летательных аппаратов важна политика freshness: если sensing-информация устарела, сервис не должен считать детекцию актуальной только потому, что ранее была достигнута высокая вероятность обнаружения. Это соответствует разделению sensing capability (способности радиозондирования) и freshness (свежести информации), используемому в ISAC-метриках \cite{liu_senscap,liu_coop}.

\section{SLA для обнаружения малоразмерных БПЛА}

Для сервиса обнаружения малоразмерных БПЛА вводится специализированный SLA-набор:
\[
\begin{array}{l}
SLA_{UAV}=(D_{admission},D_{event},D_{update},D_{fresh},\\
\qquad P_D^{min},P_{FA}^{max},P_{Miss}^{max},Acc^{min},Cov^{min}).
\end{array}
\]

В обычном timed automaton эти численные параметры не проверяются как вероятностные оценки. Они используются estimator-слоем и SDN/RIC для формирования конечных классов:
\[
\begin{array}{l}
\code{DetectionFreshnessClass},\quad \code{UpdatePeriodClass},\\
\code{FalseAlarmClass},\quad \code{MissedDetectionClass},\\
\code{AccuracyClass},\quad \code{CoverageClass}.
\end{array}
\]

Сервис считает sensing-SLA выполненным только при одновременном выполнении условий:
\[
\begin{array}{l}
SLA^{sens}_{OK}\equiv
\code{DetectionFreshnessClass=DET_FRESH}\land \code{UpdatePeriodClass=UPD_OK}\\
\qquad {}\land \code{FalseAlarmClass=FA_OK}\land
\code{MissedDetectionClass=MISS_OK}\\
\qquad {}\land \code{AccuracyClass=OK}\land
\code{CoverageClass=OK}.
\end{array}
\]

Предупреждение SLA:
\[
\begin{array}{l}
SLA^{sens}_{WARN}\equiv
\code{DetectionFreshnessClass=DET_STALE}\vee \code{UpdatePeriodClass=UPD_SLOW}\\
\qquad {}\vee \code{FalseAlarmClass=FA_HIGH}\vee
\code{MissedDetectionClass=MISS_HIGH}.
\end{array}
\]

Нарушение SLA:
\[
\begin{array}{l}
SLA^{sens}_{VIOLATED}\equiv
\code{DetectionFreshnessClass=DET_EXPIRED}\vee \code{UpdatePeriodClass=UPD_VIOLATED}\\
\qquad {}\vee \code{FalseAlarmClass=FA_CRITICAL}\vee
\code{MissedDetectionClass=MISS_CRITICAL}\\
\qquad {}\vee \code{AccuracyClass=FAILED}\vee
\code{CoverageClass=FAILED}.
\end{array}
\]

Degraded service (деградированный сервис) допустим только если причина деградации явно указана и не противоречит критичности:
\[
\begin{array}{l}
\code{AcceptedDegraded}\Rightarrow
\code{DegradedReason}\ne \code{DEG_NONE},\\
\code{CRIT_SAFETY}\land
(\code{MISS_CRITICAL}\vee \code{DetectionFreshnessClass=DET_EXPIRED})\\
\qquad {}\Rightarrow \neg\code{AcceptedDegraded}.
\end{array}
\]

Смысл этого ограничения: для safety-critical detection (критичного обнаружения) нельзя считать сервис допустимо деградированным, если высокая вероятность пропуска цели или просроченная sensing-информация делают обнаружение физически неактуальным.

\section{UPPAAL-спецификация сервисного уровня}

В UPPAAL сервисный уровень работает с конечными SLA-классами и общими переменными, установленными SDN/RIC до синхронизации. Каналы \code{service_accept?}, \code{service_degraded?}, \code{service_reject?} и \code{service_kpi_update?} не передают payload; решение допуска, причина деградации и detection KPI читаются из shared variables.

\subsection{Declarations сервисного уровня}

\begin{verbatim}
// ----- Service constants -----
const int D_req              = 3;
const int D_admission        = 15;
const int D_sla_warn         = 5;
const int D_sla_violation    = 3;
const int D_safety           = 2;
const int D_fresh_report     = 3;
const int D_update_report    = 3;
const int D_fa_report        = 3;
const int D_miss_report      = 2;
const int D_detect_warn      = 5;
const int D_detect_violation = 3;
const int D_sensing_fresh    = 10;

// ----- Service clocks -----
clock c_req, c_admission, c_sla_warn, c_sla_violation;
clock c_session, age_sensing, c_detect;

// ----- Service channels -----
broadcast chan new_demand, service_request;
broadcast chan service_accept, service_degraded, service_reject;
broadcast chan service_kpi_update, degraded_notice, freshness_expired;
broadcast chan sla_warn, sla_violation, service_complete;
broadcast chan sla_warning_report, sla_violation_report;
broadcast chan service_reconfigure, degraded_accept_report, service_terminate;

// ----- Service classes -----
typedef int[0,5] ServiceClass_t;
const ServiceClass_t SVC_URLLC=0, SVC_EMBB=1, SVC_MMTC=2;
const ServiceClass_t SVC_ISAC_SENSING=3, SVC_UAV=4, SVC_INDUSTRIAL=5;

typedef int[0,3] CriticalityClass_t;
const CriticalityClass_t CRIT_LOW=0, CRIT_MED=1, CRIT_HIGH=2, CRIT_SAFETY=3;

typedef int[0,2] SLAClass_t;
const SLAClass_t SLA_OK=0, SLA_WARN=1, SLA_VIOLATED=2;

typedef int[0,3] AdmissionClass_t;
const AdmissionClass_t ADM_PENDING=0, ADM_ACCEPTED=1;
const AdmissionClass_t ADM_REJECTED=2, ADM_DEGRADED=3;

typedef int[0,2] DemandClass_t;
const DemandClass_t DEMAND_COMM=0, DEMAND_SENS=1, DEMAND_JOINT=2;

typedef int[0,2] DetectionFreshnessClass_t;
const DetectionFreshnessClass_t DET_FRESH=0, DET_STALE=1, DET_EXPIRED=2;

typedef int[0,2] UpdatePeriodClass_t;
const UpdatePeriodClass_t UPD_OK=0, UPD_SLOW=1, UPD_VIOLATED=2;

typedef int[0,2] FalseAlarmClass_t;
const FalseAlarmClass_t FA_OK=0, FA_HIGH=1, FA_CRITICAL=2;

typedef int[0,2] MissedDetectionClass_t;
const MissedDetectionClass_t MISS_OK=0, MISS_HIGH=1, MISS_CRITICAL=2;

typedef int[0,2] QualityClass_t;
const QualityClass_t OK=0, LIMITED=1, FAILED=2;

typedef int[0,5] DegradedReason_t;
const DegradedReason_t DEG_NONE=0, DEG_STALE=1, DEG_LOW_PD=2;
const DegradedReason_t DEG_HIGH_PFA=3, DEG_UPDATE_SLOW=4;
const DegradedReason_t DEG_RESOURCE_LIMITED=5;

// ----- shared variables: payload replacement -----
ServiceClass_t serviceClass;
CriticalityClass_t criticalityClass;
SLAClass_t slaClass;
AdmissionClass_t admissionClass;
DemandClass_t demandClass;
DetectionFreshnessClass_t detectionFreshnessClass;
UpdatePeriodClass_t updatePeriodClass;
FalseAlarmClass_t falseAlarmClass;
MissedDetectionClass_t missedDetectionClass;
QualityClass_t accuracyClass;
QualityClass_t coverageClass;
DegradedReason_t degradedReason;

bool service_request_pending = false;
bool sla_warning_report_sent = false;
bool sla_violation_report_sent = false;
bool degraded_allowed = false;
bool degraded_accept_pending = false;
bool violation_report_sent = false;
\end{verbatim}

\subsection{Формулы сервисных метрик и SLA-классификации}

Целочисленная оценка риска сервиса:
\[
\begin{array}{l}
Risk_{SVC}=3M_{miss}+2M_{fa}+2M_{fresh}+M_{upd}
+M_{acc}+M_{cov}+M_{crit},
\end{array}
\]
где \(M_{miss}\), \(M_{fa}\), \(M_{fresh}\), \(M_{upd}\), \(M_{acc}\), \(M_{cov}\), \(M_{crit}\) --- численные коды соответствующих классов. В UPPAAL:
\begin{verbatim}
int svc_risk() {
  return 3*missedDetectionClass + 2*falseAlarmClass
       + 2*detectionFreshnessClass + updatePeriodClass
       + accuracyClass + coverageClass + criticalityClass;
}
\end{verbatim}

Функции SLA:
\begin{verbatim}
bool sla_sens_ok() {
  return detectionFreshnessClass == DET_FRESH
      && updatePeriodClass == UPD_OK
      && falseAlarmClass == FA_OK
      && missedDetectionClass == MISS_OK
      && accuracyClass == OK
      && coverageClass == OK;
}

bool sla_sens_warn() {
  return detectionFreshnessClass == DET_STALE
      || updatePeriodClass == UPD_SLOW
      || falseAlarmClass == FA_HIGH
      || missedDetectionClass == MISS_HIGH
      || accuracyClass == LIMITED
      || coverageClass == LIMITED;
}

bool sla_sens_violated() {
  return detectionFreshnessClass == DET_EXPIRED
      || updatePeriodClass == UPD_VIOLATED
      || falseAlarmClass == FA_CRITICAL
      || missedDetectionClass == MISS_CRITICAL
      || accuracyClass == FAILED
      || coverageClass == FAILED;
}
\end{verbatim}

\subsection{Полные guards сервисных политик}

\begin{verbatim}
bool gAdmissionRequest() {
  return serviceClass == SVC_UAV
      || serviceClass == SVC_ISAC_SENSING
      || demandClass == DEMAND_JOINT
      || demandClass == DEMAND_SENS;
}

bool gAcceptDegraded() {
  return admissionClass == ADM_DEGRADED
      && degraded_allowed
      && degradedReason != DEG_NONE
      && !(criticalityClass == CRIT_SAFETY
           && (missedDetectionClass == MISS_CRITICAL
               || detectionFreshnessClass == DET_EXPIRED));
}

bool gRejectDegraded() {
  return admissionClass == ADM_DEGRADED && !gAcceptDegraded();
}

bool gSlaWarning() {
  return slaClass == SLA_WARN || sla_sens_warn();
}

bool gSlaViolation() {
  return slaClass == SLA_VIOLATED || sla_sens_violated();
}

bool gSafetyCriticalViolation() {
  return criticalityClass == CRIT_SAFETY && gSlaViolation();
}
\end{verbatim}

\subsection{UPPAAL templates сервисного уровня}

\begin{longtable}{p{0.18\linewidth}p{0.28\linewidth}p{0.24\linewidth}p{0.22\linewidth}}
\hline
\textbf{Template} & \textbf{Locations} & \textbf{Invariants} & \textbf{Ключевые edges}\\
\hline
\(\A{REQ}\) & \code{ServiceIdle}, \code{RequestBuild}, \code{RequestPending}, \code{Accepted}, \code{AcceptedDegraded}, \code{Rejected}, \code{Completed} & \code{RequestBuild}: \(c_{req}\le D_{req}\); \code{RequestPending}: \(c_{admission}\le D_{admission}\). & \code{new_demand?}; set request variables; \code{service_request!}; accept/degraded/reject.\\
\(\A{SLA}\) & \code{SLAOk}, \code{SLAWarning}, \code{SLAViolated}, \code{SLAReported} & \code{SLAWarning}: \(c_{sla_warn}\le D_{sla_warn}\); \code{SLAViolated}: \(c_{sla_violation}\le D_{sla_violation}\). & \code{service_kpi_update?}; warning report; violation report.\\
\(\A{CRIT}\) & \code{NonCritical}, \code{Critical}, \code{SafetyCritical} & Нет отдельных часов. & update criticality; safety response guard.\\
\(\A{SVC\_AGG}\) & \code{DemandBuild}, \code{DemandSent}, \code{AwaitDecision}, \code{DecisionReceived} & \code{AwaitDecision}: \(c_{admission}\le D_{admission}\). & shared request write; \code{service_request!}; decision read.\\
\hline
\end{longtable}

Типовой edge для service request:
\begin{verbatim}
ServiceIdle -> RequestBuild
  sync new_demand?
  update c_req = 0,
         service_request_pending = false;

RequestBuild -> RequestPending
  guard gAdmissionRequest()
  sync service_request!
  update admissionClass = ADM_PENDING,
         service_request_pending = true,
         c_admission = 0;

RequestPending -> Accepted
  guard admissionClass == ADM_ACCEPTED
  sync service_accept?
  update service_request_pending = false,
         age_sensing = 0;

RequestPending -> AcceptedDegraded
  guard gAcceptDegraded()
  sync service_degraded?
  update service_request_pending = false,
         degraded_accept_pending = true;

RequestPending -> Rejected
  guard admissionClass == ADM_REJECTED || gRejectDegraded()
  sync service_reject?
  update service_request_pending = false;
\end{verbatim}

Отчет о принятии деградированного режима реализуется отдельным edge:
\begin{verbatim}
AcceptedDegraded -> AcceptedDegraded
  guard degraded_accept_pending
  sync degraded_accept_report!
  update degraded_accept_pending = false,
         violation_report_sent = false;
\end{verbatim}

\subsection{Observer-автоматы сервисного уровня}

\begin{verbatim}
template ObsServiceAdmission() {
  clock x;
  state Idle, Waiting, Violation;
  init Idle;
  Idle -> Waiting guard service_request_pending update x = 0;
  Waiting -> Idle
    guard admissionClass == ADM_ACCEPTED
       || admissionClass == ADM_DEGRADED
       || admissionClass == ADM_REJECTED;
  Waiting -> Violation guard x > D_admission;
}

template ObsDetectionFreshness() {
  clock x;
  state Idle, Waiting, Violation;
  init Idle;
  Idle -> Waiting
    guard demandClass == DEMAND_SENS
       && detectionFreshnessClass == DET_EXPIRED
    update x = 0;
  Waiting -> Idle
    guard sla_warning_report_sent || sla_violation_report_sent;
  Waiting -> Violation guard x > D_fresh_report;
}

template ObsFalseAlarmCritical() {
  clock x;
  state Idle, Waiting, Violation;
  init Idle;
  Idle -> Waiting guard falseAlarmClass == FA_CRITICAL update x = 0;
  Waiting -> Idle guard sla_violation_report_sent;
  Waiting -> Violation guard x > D_fa_report;
}

template ObsMissedDetectionCritical() {
  clock x;
  state Idle, Waiting, Violation;
  init Idle;
  Idle -> Waiting guard missedDetectionClass == MISS_CRITICAL update x = 0;
  Waiting -> Idle
    guard sla_violation_report_sent
       || violation_report_sent;
  Waiting -> Violation guard x > D_miss_report;
}
\end{verbatim}

Минимальные UPPAAL queries:
\begin{verbatim}
A[] not deadlock
A[] not ObsServiceAdmission.Violation
A[] not ObsDetectionFreshness.Violation
A[] not ObsFalseAlarmCritical.Violation
A[] not ObsMissedDetectionCritical.Violation
A[] not (criticalityClass == CRIT_SAFETY
         && missedDetectionClass == MISS_CRITICAL
         && admissionClass == ADM_DEGRADED)
\end{verbatim}

\section{Operational semantics}

\subsection{Часы и инварианты}

\begin{longtable}{p{0.24\linewidth}p{0.66\linewidth}}
\hline
\textbf{Часы} & \textbf{Назначение}\\
\hline
\(c_{req}\) & Время формирования сервисного запроса.\\
\(c_{admission}\) & Время ожидания решения SDN/RIC.\\
\(c_{sla\_warn}\) & Время реакции на предупреждение SLA.\\
\(c_{sla\_violation}\) & Время реакции на нарушение SLA.\\
\(c_{session}\) & Время активной сервисной сессии.\\
\(age_{sensing}\) & Возраст последней sensing-информации для сервиса.\\
\hline
\end{longtable}

\[
\begin{array}{ll}
\code{RequestPending}: & c_{admission}\le D_{admission},\\
\code{SLAWarning}: & c_{sla\_warn}\le D_{sla\_warn},\\
\code{SLAViolated}: & c_{sla\_violation}\le D_{sla\_violation},\\
\code{Accepted}: & age_{sensing}\le D_{sensing\_fresh}\quad \text{if sensing demand is active}.
\end{array}
\]

\subsection{Каналы}

Входные каналы:
\[
\{\code{new_demand?},\code{service_accept?},\code{service_degraded?},
\code{service_reject?},\code{sla_warn?},\code{sla_violation?},
\code{service_kpi_update?},\code{degraded_notice?},
\code{freshness_expired?},\code{service_complete?}\}.
\]

Выходные каналы:
\[
\{\code{service_request!},\code{sla_warning_report!},
\code{sla_violation_report!},\code{service_reconfigure!},
\code{degraded_accept_report!},\code{service_terminate!}\}.
\]

\section{Контракт сервисного уровня}

\[
\Contract{SVC}=(Ass_{SVC},Gar_{SVC}).
\]

\subsection{Предположения}

\begin{itemize}
  \item SDN/RIC возвращает \code{service_accept!}, \code{service_degraded!} или \code{service_reject!} до \(D_{admission}\).
  \item SLA-наблюдения поступают в конечной форме \code{SLA_OK}, \code{SLA_WARN} или \code{SLA_VIOLATED}.
  \item Сервисный запрос задает конечный класс сервиса, критичности и требований.
  \item SDN/RIC передает detection KPI (показатели обнаружения) только в дискретных классах: \code{DetectionFreshnessClass}, \code{FalseAlarmClass}, \code{MissedDetectionClass}, \code{UpdatePeriodClass}.
  \item Нижние уровни не передают сырые физические величины напрямую в сервисный автомат.
\end{itemize}

\subsection{Гарантии}

\begin{itemize}
  \item Сервис не формирует запрос вне допустимого множества классов.
  \item Если запрос принят, сервис переходит в \code{Accepted} или \code{AcceptedDegraded} и начинает контроль SLA.
  \item Если \code{SLAClass=SLA_WARN}, формируется \code{sla_warning_report!} до \(D_{sla\_warn}\).
  \item Если \code{SLAClass=SLA_VIOLATED}, формируется \code{sla_violation_report!} до \(D_{sla\_violation}\).
  \item Если \code{FalseAlarmClass=FA_HIGH} или \code{MissedDetectionClass=MISS_HIGH}, сервис обязан сформировать warning-отчет или запрос реконфигурации до \(D_{detect\_warn}\).
  \item Если \code{FalseAlarmClass=FA_CRITICAL}, \code{MissedDetectionClass=MISS_CRITICAL} или \code{DetectionFreshnessClass=DET_EXPIRED}, сервис обязан сформировать violation-отчет до \(D_{detect\_violation}\).
  \item Если запрос отклонен, \code{Rejected} содержит причину: \code{policy_reject}, \code{resource_unavailable}, \code{stale_telemetry}, \code{recovery_failed} или \code{deadline_impossible}.
\end{itemize}

\section{Переходы \(\A{REQ}\) и \(\A{SLA}\)}

Жизненный цикл запроса:
\[
\begin{array}{rcl}
\code{ServiceIdle} & \xrightarrow{\code{new_demand?}} & \code{RequestBuild},\\
\code{RequestBuild} & \xrightarrow{\code{service_request!}} & \code{RequestPending},\\
\code{RequestPending} & \xrightarrow{\code{service_accept?}} & \code{Accepted},\\
\code{RequestPending} & \xrightarrow{\code{service_degraded?}} & \code{AcceptedDegraded},\\
\code{RequestPending} & \xrightarrow{\code{service_reject?}} & \code{Rejected},\\
\code{RequestPending} & \xrightarrow{\code{timeout}} & \code{Rejected},\\
\code{Accepted} & \xrightarrow{\code{service_complete?}} & \code{Completed}.
\end{array}
\]

Деградированный режим после принятия сервиса:
\[
\begin{array}{rcl}
\code{Accepted} & \xrightarrow{\code{degraded_notice?}} & \code{AcceptedDegraded},\\
\code{AcceptedDegraded} & \xrightarrow{\code{service_kpi_update?}\land SLA^{sens}_{OK}} & \code{Accepted},\\
\code{AcceptedDegraded} & \xrightarrow{\code{sla_violation?}} & \code{SLAViolated},\\
\code{AcceptedDegraded} & \xrightarrow{\code{freshness_expired?}} & \code{SLAViolated}.
\end{array}
\]

Контроль SLA:
\[
\begin{array}{rcl}
\code{SLAOk} & \xrightarrow{\code{sla_warn?}} & \code{SLAWarning},\\
\code{SLAWarning} & \xrightarrow{\code{sla_warning_report!}} & \code{SLAOk}\ \text{or}\ \code{SLAViolated},\\
\code{SLAWarning} & \xrightarrow{\code{sla_violation?}} & \code{SLAViolated},\\
\code{SLAViolated} & \xrightarrow{\code{sla_violation_report!}} & \code{SLAReported},\\
\code{SLAReported} & \xrightarrow{\code{service_reconfigure?}} & \code{SLAOk}\ \text{or}\ \code{ServiceTerminate}.
\end{array}
\]

\section{Интерфейс требований и отчетов}

Минимальный сервисный запрос:
\[
\begin{array}{l}
R_{SVC}=(\code{ServiceClass},\code{CriticalityClass},\code{DemandClass},\\
\qquad \code{LatencyBound},\code{ReliabilityClass},\code{SensingUpdateBound},\\
\qquad \code{SensingFreshnessBound},\code{SensingAccuracyClass},\\
\qquad \code{CoverageClass},\code{PdMinClass},\code{PfaMaxClass},\\
\qquad \code{PmissMaxClass}).
\end{array}
\]

Минимальный отчет о нарушении:
\[
\begin{array}{l}
V_{SLA}=(\code{SLAClass},\code{ViolationReason},\code{ObservedAge},\\
\qquad \code{FalseAlarmClass},\code{MissedDetectionClass},\\
\qquad \code{DegradedReason},\code{DecisionRequired}).
\end{array}
\]

Для сервиса обнаружения малоразмерных БПЛА критичны не только задержка доставки события, но и давность sensing-информации. Устаревшая информация может быть формально доставлена вовремя, но уже не соответствовать текущему положению цели.

\section{Верификационные свойства и observer-автоматы}

\[
\begin{array}{l}
A[]\ \neg\code{deadlock},\\
A[]\ (\code{RequestPending}\Rightarrow c_{admission}\le D_{admission}),\\
A[]\ (\code{RequestPending}\Rightarrow \Diamond(\code{Accepted}\vee \code{AcceptedDegraded}\vee \code{Rejected})),\\
A[]\ (\code{SLAViolated}\Rightarrow \Diamond \code{sla_violation_report_sent}),\\
A[]\ \neg(\code{Accepted}\land \code{reason=resource_unavailable}).
\end{array}
\]

Для safety-critical сервиса:
\[
\code{CRIT_SAFETY}\land \code{SLA_VIOLATED}\Rightarrow
\Diamond_{\le D_{safety}}
(\code{service_reconfigure}\vee \code{service_terminate}\vee \code{violation_report_sent}).
\]

Для freshness:
\[
\code{DemandClass=DEMAND_SENS}\land \code{DetectionFreshnessClass=DET_EXPIRED}
\Rightarrow
\Diamond_{\le D_{fresh\_report}}\code{sla_warning_report_sent}.
\]

Для периода обновления sensing:
\[
\code{UpdatePeriodClass=UPD_VIOLATED}\Rightarrow
\Diamond_{\le D_{update\_report}}\code{sla_violation_report_sent}.
\]

Для риска ложной тревоги:
\[
\code{FalseAlarmClass=FA_CRITICAL}\Rightarrow
\Diamond_{\le D_{fa\_report}}\code{sla_violation_report_sent}.
\]

Для риска пропуска цели:
\[
\code{MissedDetectionClass=MISS_CRITICAL}\Rightarrow
\Diamond_{\le D_{miss\_report}}
(\code{service_reconfigure}\vee \code{service_terminate}\vee \code{sla_violation_report_sent}).
\]

\section{Ограничения применимости}

\begin{itemize}
  \item SLA должен быть переведен в конечные классы; текстовое SLA без \(\alpha_{SVC}\) не проверяется timed automata.
  \item Сервисный уровень не доказывает физическую вероятность обнаружения и не вычисляет CRB.
  \item AcceptedDegraded допустим только при явном сервисном разрешении деградированного режима.
  \item Для БПЛА-сценария detection freshness является отдельным ограничением, а не синонимом вероятности обнаружения.
  \item False alarm и missed detection задаются как дискретные классы внешней оценки; сервисный автомат не оценивает статистику обнаружения по выборке.
  \item Если service layer напрямую управляет PHY/MAC, иерархическая формализация нарушается.
\end{itemize}

\section{Заключение}

Application / Service Layer задается как композиция автоматов запроса, SLA, критичности и агрегирования требований. Его задача --- формализовать сервисные ограничения и проверять своевременное получение решения или отчета о нарушении. Слабое место ранней модели состоит в том, что SLA может остаться текстовым описанием без конечных классов. В данной спецификации это закрыто отображением \(\alpha_{SVC}\), явными классами \code{AdmissionClass}, \code{SLAClass}, \code{CriticalityClass}, \code{DetectionFreshnessClass}, \code{FalseAlarmClass}, \code{MissedDetectionClass} и observer-свойствами для admission, SLA violation, detection freshness, update period, false alarm и missed detection.

\begin{thebibliography}{9}

\bibitem{liu_senscap}
G. Liu et al., ``SensCAP: A Systematic Sensing Capability Performance Metric for 6G ISAC,'' \emph{IEEE Internet of Things Journal}, vol. 11, no. 18, pp. 29438--29454, 2024, doi: 10.1109/JIOT.2024.3430502.

\bibitem{liu_coop}
G. Liu et al., ``Cooperative Sensing for 6G ISAC: Concept, Key Technologies, Performance Evaluation, and Field Trial,'' \emph{Engineering}, vol. 56, pp. 130--148, 2026.

\bibitem{alur_dill}
R. Alur and D. L. Dill, ``A Theory of Timed Automata,'' \emph{Theoretical Computer Science}, vol. 126, no. 2, pp. 183--235, 1994.

\bibitem{mosudi_sdn}
O. A. Mosudi, S. I. Popoola, A. A. Atayero, and A. A. Elngar, ``SDN for 5G and Beyond: Transforming Network Architecture,'' 2025.

\bibitem{zolanvari_sdn}
Zolanvari, ``SDN for 5G,'' 2015.

\end{thebibliography}

\end{document}
